\section{Quantum Programming Languages}\label{sec:lit_rev:qpl}
Programming languages are the tool with which computer scientists and software engineers convey their algorithmic intent to machines for execution.
As described in Section~\ref{sec:lit_rev:abstraction} programming languages come in many types and families, each serving a different role within the execution stack or compilation chain.
These roles include high-level languages, designed to make human intent as easy as possible to express, intermediate representations which make program structure explicit for compiler analysis and transformation and assembly languages which describe execution in terms machines directly operate on.

This section explores the current state of quantum programming languages; those intended to execute on quantum computers.
It examines which layers of the compilation chain are represented, what they offer and whether they support the expression of human intent without requiring programmers to reason at the level of the quantum machine.
The evidence shows genuine progress across all layers; no existing system, however, enables sustained human-centred reasoning: the programmer's required mental model remains machine-centred throughout.

\subsection{Scope and Vocabulary}\label{subsec:lit_rev:qpl:scope}
The review presented in this section will analyse quantum programming languages across four categories or families.
Each of those families is described in Table~\ref{tab:lit_rev:qpl:language_families}.
The definitions of these families, their taxonomy and critical analyses are expanded in Subsection~\ref{subsec:lit_rev:qpl:taxonomy}.
Other frameworks assess quantum languages on their semantic or syntactic principles, such as imperative and functional languages~\cite{Gay2006QPLSurvey}.
The review presented here does not do that and instead focuses on the intended purpose in the compilation stack.

\begin{table}[htbp]
    \centering
    \caption{Summary of quantum programming languages families identified in the current literature.}
    \label{tab:lit_rev:qpl:language_families}
    \begin{tabularx}{\textwidth}{@{}lX@{}}
        \toprule
        \textbf{Family} & \textbf{Description}          \\ \midrule
        High-level & Human-readable algorithmic logic; quantum resource management abstractions. \\
        Embedded framework & Frameworks or libraries for interacting with quantum hardware hosted in classical languages like Python, also referred to as hosted languages~\cite{corralesgarro2025productivityexpressiveness}. \\
        Intermediate representation (IR) & Machine-independent optimization and canonical representation (IR). \\
        Assembly & Mapping logic to hardware-level gate sequences and qubit registers. \\ \bottomrule
    \end{tabularx}
\end{table}

Each family represents a different role in the compilation and execution of quantum programs or provides ergonomic or abstraction benefits.
The families of languages selected represent languages that are used by programmers in the development of quantum programs or those required to demonstrate the full compilation chains and comparisons of complexity.
From those criteria languages eligible for inclusion are embedded frameworks in classical languages, intermediate representation languages, quantum assembly languages and high-level quantum development languages.
Languages that solely represent hardware control and vendor or specific machine level languages are excluded from review as they are either not intended for human development or fail to add value to discussions about abstraction.

To evaluate the current state of quantum programming languages, the review will consider the following elements of each language and family:
\begin{itemize}
    \item Abstraction level: The degree to which a system exposes machine-level quantum detail to a programmer~\cite{DiMatteo2024AbstractionHierarchy}.
    \item Required quantum knowledge: What quantum mechanical understanding a programmer must have to productively use the language~\cite{Myers2016PLUsabilitySIG, ferreira2024qplusage}.
    \item System portability: Whether the language can be used to target multiple hardware backends without programmer intervention~\cite{DiMatteo2024AbstractionHierarchy, nunez_corrales2025betterabstractmachines}.
\end{itemize}
Together these criteria represent existing evaluation frameworks and known areas of friction for practitioners of quantum programming.
The level to which a language adheres to them is a signifier of its ability to advance the current state of quantum programming languages.

Each language family will have a different aspirational level for the criteria presented, intermediate representations will require a higher degree of quantum knowledge than high-level languages.
The criteria are applied contextually based on the purpose of each language family.

\subsection{Taxonomy and the Missing Tier}\label{subsec:lit_rev:qpl:taxonomy}
Sections~\ref{subsec:lit_rev:abstraction:levels_elements} and~\ref{subsec:lit_rev:abstraction:looking_forward} established that quantum programming lacks a machine-independent source layer.
This section grounds that finding in an examination of what each language family aspires to and whether any family achieves it.

The defining aspiration of quantum programming language development is quantum-agnostic application development~\cite{nunez_corrales2025betterabstractmachines}.
A family of languages that fulfills this aspiration would not require fundamental knowledge of quantum mechanics or quantum circuits to program.
Variations of this aspiration exist; some argue that quantum agnosticism could be achieved by improving circuit generation techniques to remove quantum pre-requisite knowledge~\cite{cobb2022higherlevelabstractions}.
Others identify a lack of separation between instructions and circuit primitives as preventing the formulation of a genuine source-layer language~\cite{nunez_corrales2025betterabstractmachines}.
Additional analysis has shown a lack of complex data structures and quantum control sequence support~\cite{Gay2006QPLSurvey} as further barriers.

Regardless of the path, the literature converges on fundamental principles such a language family would adhere to:
\begin{itemize}
    \item Minimal or no understanding of fundamental quantum mechanics required to write programs.
    \item No direct manipulation of the quantum circuit model or equivalents in other hardware.
    \item The programmers mental model exists in the problem space; not quantum mechanics or circuits.
\end{itemize}

Four language families are identified in the literature and introduced in Subsection~\ref{subsec:lit_rev:qpl:scope}:
\begin{itemize}
    \item High-level
    \item Embedded framework
    \item Intermediate representation
    \item Assembly
\end{itemize}

Each offers different aspirational goals in the quantum programming domain.
A summary is provided in Table~\ref{tab:lit_rev:qpl:language_family_taxonomy}.
This subsection characterises each family; Subsection~\ref{subsec:lit_rev:qpl:reading} examines them in detail.
\begin{table}[htbp]
    \centering
    \caption{Summary of quantum programming languages families and their defined aspirations.}
    \label{tab:lit_rev:qpl:language_family_taxonomy}
    \begin{tabularx}{\textwidth}{@{}>{\hsize=0.2\hsize}X>{\hsize=0.2\hsize}X>{\hsize=0.6\hsize}X@{}}
        \toprule
        \textbf{Family} & \textbf{Language Example} & \textbf{Description}          \\ \midrule
        High-level & Silq & Algorithmic expression with compiler-managed quantum obligations. \\
        Embedded framework & Qiskit & Ecosystem integration, hardware portability and simulation access; hosted in familiar classical environments.  \\
        Intermediate representation (IR) & QIR & Stable compilation substrate enabling optimisation passes, toolchain portability across front-ends, and a formal verification surface for cross-layer equivalence. \\
        Assembly & OpenQASM & Portable, complete execution-level specification; hardware-independent circuit interchange. \\ \bottomrule
    \end{tabularx}
\end{table}

Assembly languages such as OpenQASM aspire to offer a portable, hardware-agnostic specification of quantum circuits.
They describe circuit primitives alongside hardware concerns such as pulses, timing delays and qubit rotations or gates~\cite{cross2022openqasm3}.
OpenQASM achieves these aspirations.

Intermediate representations (IRs) offer the same aspirations for quantum programming as their classical counterparts.
A productive IR is machine-independent and stable across the compilation stack; both forward and backward compatible~\cite{irreview2024}.
IRs retain the intent of the source layer while enabling compiler optimisation passes and formal verification~\cite{irreview2024}.
Current quantum IRs may or may not encode quantum-physical constraints such as no-cloning or teleportation.
QIR and Q-MLIR achieve these outcomes~\cite{irreview2024}.

Embedded frameworks enable ecosystem integration for remote execution and simulation of quantum programs from within a classical programming language, providing ergonomic tools for common quantum algorithms and circuit compilation for quantum execution systems.
Qiskit~\cite{aleksandrowicz2019qiskit} is a representative example, offering Python libraries for circuit construction, compilation across multiple hardware backends and remote execution.
PennyLane~\cite{bergholm2018pennylane} offers comparable tooling with a focus on variational quantum circuits and machine learning techniques.

High-level languages aspire to provide a layer where algorithmic thinking drives the programmer.
These languages aim to provide an abstraction level where the compiler can reason about quantum obligations such as physical correctness and resource management.
As Section~\ref{subsec:lit_rev:abstraction:in_quantum_computing} showed, Silq, Qmod and Q\# offer many high-level tools to support algorithmic programming.
However, many of these tools still require the programmer to understand the quantum-mechanical and circuit-level principles underlying their programs.
High-level languages do not divorce the quantum from quantum programming.

The taxonomy reveals a missing ``fifth-tier'' of quantum programming languages.
The quantum-agnostic, application programming paradigm aspired to in the literature is not realisable with any of the language families presented.
Assembly and intermediate representations are explicitly quantum by design, requiring a fundamental understanding of circuit primitives and hardware concerns.
Embedded frameworks were not intended to produce quantum-agnosticism; their goal is to improve access to near-term quantum computers and remote execution of quantum circuit definitions.
Embedded frameworks achieve this effectively, but it is not an adequate solution to the missing tier.
High-level languages come closest, but their fundamental coupling to quantum circuit primitives and resource management requires the programmer to have a quantum mechanical mental model.
These languages do not enable problems and solutions to be formulated purely algorithmically.

\subsection{Reading the Language Landscape}\label{subsec:lit_rev:qpl:reading}
Quantum programming language development continues to focus on improving computability and tractability within quantum domains rather than focusing on improving the experience of programmers~\cite{corralesgarro2025productivityexpressiveness}.
Assessing programming languages solely against their computability biases success towards machine concerns rather than human factors and programmer experience.

Current quantum programming languages suffer from a usability issue failing to be addressed at a fundamental level~\cite{ferreira2024qplusage, nunez_corrales2025betterabstractmachines}.
A survey of active quantum programming projects on GitHub showed a lack of novel development or algorithmic investigation~\cite{murillo2022sequantumprogramming}.
Surveys of embedded framework usage found that their popularity can be attributed to the improved programmer experience when using a classical language; however, programmers still struggled with the circuit architecture~\cite{ferreira2024qplusage}.
Programmers consistently report quantum domain knowledge, vocabulary and concepts as barriers to entry and understanding of quantum languages~\cite{ferreira2024qplusage, murillo2022sequantumprogramming}.

Subsection~\ref{subsec:lit_rev:qpl:scope} established three criteria that can be used to critically examine quantum programming languages, these criteria are directly related to the programmer experience:
\begin{itemize}
    \item Abstraction level.
    \item Required quantum knowledge.
    \item System portability.
\end{itemize}

High-level programming languages realise the abstraction requirements described in Section~\ref{sec:lit_rev:abstraction}.
In classical programming language theory they provide a control plane that reduces a programmer's exposure to machine concerns~\cite{scott2009plp}.

In the remainder of this subsection, the various language families in quantum computing are described in detail.
Their definitions are presented and members of the family analysed for their features and success as programming languages.
Languages are assessed on the degree to which they support computational thinking paradigms~\cite{Wing2006ComputationalThinking}, graded based on their family and position in the compilation stack.

Each language family is assessed on its suitability to support the ideal quantum abstract machine~\cite{nunez_corrales2025betterabstractmachines} described in Section~\ref{sec:lit_rev:abstraction} with respect to its role in compilation.
High-level languages are assessed with additional rigour against their ability to shield programmers from inherent constraints in quantum programming not present in classical programming.
\textit{Representation exposure}; the degree to which a programmer is required to reason about internal quantum state directly, due to the inability to copy or inspect it without disturbing the computation.
\textit{Semantic mismatch}; the degree to which invalid equivalence with classical concepts is assumed~\cite{yuan2025foundationalabstractions}.

%% ------------------- TO BE WRITTEN ------------------- %%
\subsubsection{Assembly Languages}\label{subsubsec:lit_rev:qpl:assembly}

\subsubsection{Intermediate Representations}\label{subsubsec:lit_rev:qpl:ir}

\subsubsection{Embedded Frameworks}\label{subsubsec:lit_rev:qpl:frameworks}

\subsubsection{High-Level Languages}\label{subsubsec:lit_rev:qpl:highlevel}

\subsection{The Usability Ceiling}\label{subsec:lit_rev:qpl:ceiling}
- gay 2006 provides great points about whats missing and why high level abstraction is critical, they may help ground discussions